% Generated draft: requires mathematical and source review.
\documentclass[11pt]{article}
\usepackage[margin=25mm]{geometry}
\usepackage{fontspec}
\setmainfont{Latin Modern Roman}
\usepackage{amsmath,amssymb,mathtools,unicode-math}
\setmathfont{Latin Modern Math}
\usepackage{enumitem,needspace,xurl,hyperref}
\setlength{\parindent}{0pt}
\setlength{\parskip}{5pt}
\setlength{\emergencystretch}{3em}
\newcommand{\N}{\mathbb{N}}
\newcommand{\Z}{\mathbb{Z}}
\newcommand{\Q}{\mathbb{Q}}
\newcommand{\R}{\mathbb{R}}
\newcommand{\C}{\mathbb{C}}
\newcommand{\F}{\mathbb{F}}
\newcommand{\A}{\mathbb{A}}
\newcommand{\PP}{\mathbb{P}}
\newcommand{\E}{\mathbb{E}}
\newcommand{\eps}{\varepsilon}
\newcommand{\dd}{\,\mathrm{d}}
\begin{document}
\section*{UnsolvedMath ID: 701233}
\textbf{Source problem code:} catalogue entry
\subsection{Definitions and mathematical statement}
Let \(p\) be a prime and let \(n\) be an integer with \(n>1\). A pro-\(p\) group is a topological group that is an inverse limit of finite \(p\)-groups. Let \(F\) be the free pro-\(p\) group of rank \(n\): it has distinguished topological generators \(x_1,\ldots,x_n\) such that, for every pro-\(p\) group \(G\) and every \((g_1,\ldots,g_n)\in G^n\), there is a unique continuous homomorphism \(F\to G\) sending \(x_i\) to \(g_i\) for each \(i\).

Write \(\operatorname{Aut}_{\mathrm{cont}}(F)\) for the group of continuous automorphisms of \(F\), equipped with its usual profinite topology. A neighborhood basis of the identity consists of the kernels of the maps
\[
\operatorname{Aut}_{\mathrm{cont}}(F)\longrightarrow\operatorname{Aut}(F/U),
\]
where \(U\) ranges over open characteristic subgroups of \(F\). This is the conventional topological interpretation of the source's notation \(\mathrm{Aut}\,F\); the supplied statement does not explicitly specify the topology.

For a profinite group \(H\) and a prime \(\ell\), its continuous \(\ell\)-cohomological dimension \(\operatorname{cd}_{\ell}(H)\) is the least nonnegative integer \(d\), if one exists, such that
\[
H^q_{\mathrm{cont}}(H,M)=0
\quad\text{for every integer }q>d
\text{ and every discrete continuous }H\text{-module }M
\text{ whose underlying abelian group is }\ell\text{-primary torsion}.
\]
If there is no such integer, set \(\operatorname{cd}_{\ell}(H)=\infty\). The full continuous cohomological dimension is \(\operatorname{cd}(H)=\sup_{\ell\text{ prime}}\operatorname{cd}_{\ell}(H)\).

\textbf{Question (catalogue entry).} For every prime \(p\) and every free pro-\(p\) group \(F\) of finite rank \(n>1\), does \(\operatorname{Aut}_{\mathrm{cont}}(F)\) possess an open subgroup of finite cohomological dimension?

\textbf{Unresolved convention.} The supplied record does not specify whether ``cohomological dimension'' means full continuous cohomological dimension or continuous \(p\)-cohomological dimension. With the definitions above, the respective precise readings ask for an open subgroup \(H\leq\operatorname{Aut}_{\mathrm{cont}}(F)\) satisfying \(\operatorname{cd}(H)<\infty\) or \(\operatorname{cd}_p(H)<\infty\). The intended convention requires source verification; neither reading is silently selected here.
\subsection{Short English statement}
For a free pro-\(p\) group on finitely many generators, with at least two generators, ask whether its continuous automorphism group has an open subgroup of finite cohomological dimension. The supplied record labels the problem open. The intended cohomological-dimension convention remains to be verified.
\subsection{Sources}
\begin{itemize}
\item Supplied local UnsolvedMath registry record: permanent ID \(701233\), exact source code \texttt{catalogue entry}, title \emph{Kourovka Notebook Problem 16.41}; upstream revision \texttt{d371ee917ac9647fbb06cd6fc529deec6f3dda82}. The record was created and updated on 2026-10-16 and supplies the status \texttt{open}; this status is not independently verified here.
\item Attributed source: \emph{The Kourovka Notebook (Unsolved Problems in Group Theory)}, 16th issue (2006), Problem 16.41; also identified in the supplied background with the 21st issue edition (2026), arXiv identifier \texttt{1401.0300v48}. These bibliographic details are supplied attribution, not a claim of source consultation.
\item The supplied collection description names E. I. Khukhro and V. D. Mazurov as editors. The proposal year is 2006. The proposer field is the malformed string \texttt{P. A. Zalesski\textbackslash u\{\textbackslash i}; its incomplete accent sequence requires verification before the name can be typeset reliably.
\end{itemize}
\end{document}
